\documentclass[10pt,a4paper]{amsart}
\usepackage[margin=19mm]{geometry}
\usepackage{amsmath,amssymb,mathtools}
\usepackage[scr=rsfs,cal=euler]{mathalfa}
\usepackage{xfrac}
\usepackage[unicode,hidelinks,bookmarks=false]{hyperref}
\newcommand{\ie}{i.e.\ }
\newcommand{\cf}{cf.\ }
\newcommand{\resp}{respectively\ }
\newcommand{\Subsection}{\subsection}
\providecommand{\nobreakdash}{\nobreak\hbox{-}}
\setlength{\emergencystretch}{2em}
\multlinegap0pt
\usepackage{bm}
\usepackage{bbm}
\usepackage{dsfont}
\usepackage{xspace}
\usepackage{cancel}
\usepackage{mathtools}
\usepackage{amsthm}
\makeatletter
\@ifundefined{theorem}{\newtheorem{theorem}{Theorem}}{}
\makeatother
\title[New reformulations for 0-1 quadratic programming problem usi]{New reformulations for 0-1 quadratic programming problem using quadratic nonconvex reformulation techniques and valid inequalities}
\author{Cheng Lu, Yu Fei, Jing Zhou, Zhibin Deng, Guangtai Qu}
\date{Source: arXiv:2602.19051v1 (CC BY 4.0)}
\begin{document}
\maketitle

\noindent\emph{Source and licence.} New reformulations for 0-1 quadratic programming problem using quadratic nonconvex reformulation techniques and valid inequalities. Version arXiv:2602.19051v1 (CC BY 4.0), distributed under the Creative Commons Attribution 4.0 International licence, as recorded in the arXiv licence metadata for this exact version and shown on the abstract page https://arxiv.org/abs/2602.19051v1. Full rights notice: licenses/arxiv-2602.19051-CC-BY-4.0.txt.\par\smallskip

\noindent\textit{Editorial scope.} Counted unit: the Theorem whose source label is \texttt{thm1} in the article (source0.tex lines 382--388, source label \texttt{thm1}), delivered together with the article's own complete original proof of it (source0.tex lines 389--482, one \texttt{proof} environment of 94 lines). No mathematics is written, repaired or re-derived: statement and proof are the article's own text line for line. Required context retained uncounted: QNRE (lines 344--351); QCRE-MCR (lines 355--364); QNR-dual1 (lines 367--379). Every definition, display and label of those lines is retained unchanged. Omitted from this package and not redistributed: 1--343, 352--354, 365--366, 380--381, 483--955. Nothing inside a retained excerpt is omitted or altered: every retained body is delivered exactly as the article's lines are written, so no omission receipt of any kind accompanies this package. Citations: the retained excerpts carry no citation command at all, so no bibliography entry of the article is transcribed and no reference is redistributed. Reference adaptation: none was needed and none was made; every reference inside the delivered unit resolves inside it, so no unresolved reference or citation is produced. Printed slips kept literal: no printed word, symbol, hypothesis or number of the counted statement or of its proof was altered. Layout and class: the article's own class is \texttt{sn-jnl} with packages including graphicx, multirow, amsmath, amssymb, amsfonts, amsthm, mathrsfs, appendix, xcolor, textcomp, manyfoot, booktabs, algorithm, algorithmicx; the wrapper is the lane's pinned \texttt{amsart} editorial wrapper at 10pt on A4, which restates the theorem-like environments the retained text needs and adds the article's own macro definitions that the retained text needs (none), with extra wrapper packages (bm, bbm, dsfont, xspace, cancel, mathtools). The delivered numbering is the unit's own. Assets: no retained span contains a figure environment, a graphics-inclusion command, a PDF-page inclusion, a table or a TikZ picture; no image file is copied into this package, the package declares no asset dependency and no image file is redistributed with it. Licence: version arXiv:2602.19051v1 of this article is distributed under the Creative Commons Attribution 4.0 International licence, as recorded in the arXiv licence metadata for this exact version and shown on the abstract page https://arxiv.org/abs/2602.19051v1. A full rights notice is delivered at licenses/arxiv-2602.19051-CC-BY-4.0.txt. Characters: every retained excerpt is pure ASCII, so the delivered engine, XeLaTeX with its default Latin Modern text font, reports no missing character.
\begin{equation}\label{QNRE}\tag{QNRE}
	\begin{aligned}
		\min ~&\frac{1}{2}x^\top(Q+2\textrm{diag}(\lambda)-2Z)x+c^\top x-\lambda^\top x+w+\sum_{t=1}^m \gamma_t \left(g_t(x)-v_t\right)\\%~ f_{\lambda,Z,\gamma}(x,w,v)\\
		\textrm{s.t.}~&x_i \in \{0,1\},~i=1\ldots,n,\\
		&w\geq x^\top Z x,\\
		& v_t= g_t(x),~v_t\leq 0,~t=1,\ldots,m.
	\end{aligned}
\end{equation}

\begin{equation}\label{QCRE-MCR}
	\begin{aligned}
		\min ~&\frac{1}{2}x^\top(Q+2\textrm{diag}(\lambda)-2Z)x+c^\top x-\lambda^\top x+w+\sum_{t=1}^m \gamma_t \left(g_t(x)-v_t\right)\\
		\textrm{s.t.}~ &w\geq Z\cdot X,\\
		&x_i^2\leq X_{ii} \leq x_i,~i=1,\ldots,n,\\
		& v_t= \frac{1}{2}A_t \cdot X+ c_t^\top x+b_t,~t=1,\ldots,m,\\
		&v_t\leq 0,~t=1,\ldots,m,\\
		&(X_{ij},x_i,x_j)\in \mathcal{M}_{ij}, ~i\neq j, ~i,j=1,\ldots,n.\\
	\end{aligned}
\end{equation}

\begin{equation}\label{QNR-dual1}
	\begin{aligned}
		\max ~&~ \sigma \\
		\mbox{s.t.} ~&\begin{bmatrix}
			2\hat{b}-2\sigma &\hat{c}^T\\
			\hat{c} ~&\hat{Q} \\
		\end{bmatrix}\succeq 0,\\
		&\hat{Q}=Q+2\textrm{diag}(\lambda)+\sum_{t=1}^m\gamma_t A_t-2M-2N+2R+2S,\\
		&\hat{c}=c-\lambda+\sum_{t=1}^m \gamma_t c_t+2Ne-Re-Se,\\
		&\hat{b}=\sum_{t=1}^m \gamma_t b_t-\sum_{i=1}^n\sum_{j=1}^n N_{ij},\\
		&\sigma\in\mathbb{R},~\lambda\in\mathbb{R}^n,~\gamma\in\mathbb{R}^m_+,~M,N,R,S\in\mathbb{X}_+^n.
	\end{aligned}
\end{equation}

\begin{theorem}\label{thm1}
	The parameter $(\lambda^\ast,Z^\ast,\gamma^\ast)$ that maximizes the McCormick relaxation bound of \eqref{QNRE} can be obtained by setting $\lambda^\ast=\lambda^\dag$, $\gamma^\ast=\gamma^\dag$, and
	\begin{equation}
		Z^\ast= M^\dag +  N^\dag -  R^\dag -  S^\dag,
	\end{equation}
	where $(\lambda^\dag,\gamma^\dag,\tau^\dag,M^\dag,N^\dag,R^\dag,S^\dag)$ denotes the optimal solution of \eqref{QNR-dual1}.
\end{theorem}

\begin{proof}
	First, we simplify \eqref{QCRE-MCR} as follows:
	\begin{equation}\label{QNRE-MCR}
		\begin{aligned}
			\min ~&~\frac{1}{2}x^\top(Q+2\textrm{diag}(\lambda)-2Z)x+c^\top x-\lambda^\top x+ Z\cdot X+\sum_{t=1}^m \gamma_t \left(g_t(x)-v_t\right)\\
			\textrm{s.t.}~ &x_i^2\leq X_{ii} \leq x_i,~i=1,\ldots,n,\\
			& v_t= \frac{1}{2}A_t \cdot X+ c_t^\top x+b_t,~t=1,\ldots,m\\
			&v_t\leq 0,~t=1,\ldots,m,\\
			&(X_{ij},x_i,x_j)\in \mathcal{M}_{ij}, ~i\neq j, ~i,j=1,\ldots,n.\\
		\end{aligned}
	\end{equation}	
	Define the Lagrangian function of \eqref{QNRE-MCR} as
	\begin{equation}
		\begin{aligned}
			&L(x,X,v,\alpha,\beta,\tau,\zeta,M,N,R,S)=\frac{1}{2}x^\top(Q+2\textrm{diag}(\lambda)-2Z)x+c^\top x-\lambda^\top x \\
			&\quad + Z\cdot X+\sum_{t=1}^m \gamma_t \left(\frac{1}{2} x^\top A_t x+c_t^\top x+b_t-v_t\right)+\sum_{i=1}^n\alpha_i(x_i^2-X_{ii})\\
			&\quad +\sum_{i=1}^n\beta_i(X_{ii}-x_i)+\sum_{t=1}^m \tau_t\left(  v_t-\frac{1}{2}A_t \cdot X- c_t^\top x-b_t \right)
			+\sum_{t=1}^m \zeta_t v_t\\
			&\quad +\sum_{i=1}^n\sum_{j=1}^n\left[ -M_{ij}X_{ij}-N_{ij}\left(X_{ij}-x_j-x_i+1\right)  \right]\\
			&\quad +\sum_{i=1}^n\sum_{j=1}^n\left[  R_{ij}(X_{ij}-x_i)+S_{ij}(X_{ij}-x_j) \right],
		\end{aligned}
	\end{equation}
	where $\alpha,\beta\in \mathbb{R}^n_+$, $\tau\in \mathbb{R}^m$ and $\zeta\in \mathbb{R}^m_+$ are the Lagrangian multipliers of the corresponding constraints in \eqref{QNRE-MCR}, and $M, N, R, S \in \mathbb{X}^n_+$, whose entries $M_{ij}, N_{ij}, R_{ij}, S_{ij}$ are the Lagrangian multipliers of the four McCormick inequalities appear in the constraint $\left(X_{ij},x_i,x_j\right)\in \mathcal{M}_{ij}$, respectively. Define the dual function of \eqref{QNRE-MCR} as
	\begin{equation}\label{dualfunction}
		p(\alpha,\beta,\tau,\zeta,M,N,R,S)=\min_{x\in\mathbb{R}^n,X\in\mathbb{S}^n,v\in\mathbb{R}^m}L(x,X,v,\alpha,\beta,\tau,\zeta,M,N,R,S),
	\end{equation}
	and define the associated dual problem of \eqref{QNRE-MCR} as
	\begin{equation}\label{QNR-dual}
		\begin{aligned}
			\max ~&p(\alpha,\beta,\tau,\zeta,M,N,R,S)\\
			\mbox{s.t.} ~~&\alpha,\beta\in \mathbb{R}^n_+,~\tau\in \mathbb{R}^m,~\zeta\in \mathbb{R}^m_+, ~M,N,R,S\in\mathbb{X}_+^n.
		\end{aligned}
	\end{equation}
	To avoid the unboundedness in minimizing the Lagrangian function in \eqref{dualfunction}, the coefficients for the matrix $X$ and vector $v$ in the Lagrangian function must be equal to zero. This leads to the following equations:
	\begin{equation}\label{eq_condition}
		\left\{
		\begin{aligned}
			&Z-\textrm{diag}(\alpha)+\textrm{diag}(\beta)-\sum_{t=1}^m\frac{1}{2} \tau_t A_t-M-N+R+S= 0,\\
			&-\gamma_t+\tau_t+\zeta_t=0,~t=1,\ldots,m.
		\end{aligned}\right.
	\end{equation}
	Using \eqref{eq_condition}, the Lagrangian function can be simplified to
	\begin{equation}
		L(x,X,v,\alpha,\beta,\tau,\zeta,M,N,R,S) = \frac{1}{2}x^T \hat{Q} x + \hat{c}^T x + \hat{b},
	\end{equation}
	where
	\begin{align}
		&\hat{Q}=Q+2\textrm{diag}(\lambda)-2Z+\sum_{t=1}^m \gamma_t A_t +2\textrm{diag}(\alpha),\\
		&\hat{c}=c-\lambda+\sum_{t=1}^m \gamma_t c_t-\beta- \sum_{t=1}^m \tau_t c_t +2Ne-Re-Se,\\
		&\hat{b}=\sum_{t=1}^m \gamma_t b_t-\sum_{t=1}^m \tau_t b_t-\sum_{i=1}^n\sum_{j=1}^n N_{ij}.
	\end{align}
	Based on the first equation in \eqref{eq_condition}, the expression for $\hat{Q}$ can equivalently be written as
	\begin{equation}
		\hat{Q}=Q+2\textrm{diag}(\lambda)+\sum_{t=1}^m\gamma_t A_t-\sum_{t=1}^m \tau_t A_t +2\textrm{diag}(\beta)-2M-2N+2R+2S.
	\end{equation}
	The dual problem can now be equivalently transformed to the following problem:
	\begin{equation}\label{QNR-dual2}
		\begin{aligned}
			\max ~&~ \sigma \\
			\mbox{s.t.} ~&~\frac{1}{2}x^T \hat{Q} x + \hat{c}^T x + \hat{b}-\sigma \geq 0,~\forall x\in\mathbb{R}^n,\\
			&\textrm{(16), (19--21)}\\
			&\sigma\in\mathbb{R},~\alpha\in \mathbb{R}^n_+,~\beta\in \mathbb{R}^n_+,~\tau\in \mathbb{R}^m,\zeta\in \mathbb{R}^m_+,~M,N,R,S\in\mathbb{X}_+^n.
		\end{aligned}
	\end{equation}
	Besides, the constraint $\frac{1}{2}x^T \hat{Q} x + \hat{c}^T x + \hat{b}-\sigma \geq 0~(\forall x\in\mathbb{R}^n)$ is equivalent to
	\begin{equation}\label{sdpconstraint}
		\begin{bmatrix}
			2\hat{b}-2\sigma &\hat{c}^T\\
			\hat{c} ~&\hat{Q} \\
		\end{bmatrix}\succeq 0
	\end{equation}
	Thus, the first constraint in \eqref{QNR-dual2} can be replaced by \eqref{sdpconstraint}.
	
	By strong duality, we have that the optimal value of \eqref{QNRE-MCR} is equal to the optimal value of \eqref{QNR-dual2}. Hence, the problem of maximizing the optimal value of \eqref{QNRE-MCR} is equivalent to the problem of maximizing the optimal value of \eqref{QNR-dual2}, which can be formulated as follows:
	\begin{equation}\label{QNR-dual3}
		\begin{aligned}
			\max ~&~ \sigma \\
			\mbox{s.t.} ~&\begin{bmatrix}
				2\hat{b}-2\sigma &\hat{c}^T\\
				\hat{c} ~&\hat{Q} \\
			\end{bmatrix}\succeq 0,\\
			&\hat{Q}=Q+2\textrm{diag}(\lambda)+\sum_{t=1}^m\gamma_t A_t-\sum_{t=1}^m \tau_t A_t +2\textrm{diag}(\beta)-2M-2N+2R+2S,\\
			&\hat{c}=c-\lambda+\sum_{t=1}^m \gamma_t c_t-\beta- \sum_{t=1}^m \tau_t c_t +2Ne-Re-Se,\\
			&\hat{b}=\sum_{t=1}^m \gamma_t b_t-\sum_{t=1}^m \tau_t b_t-\sum_{i=1}^n\sum_{j=1}^n N_{ij},\\
			&Z-\textrm{diag}(\alpha)+\textrm{diag}(\beta)-\sum_{t=1}^m\frac{1}{2} \tau_t A_t-M-N+R+S= 0,\\
			&-\gamma_t+\tau_t+\zeta_t=0,~t=1,\ldots,m,\\
			&Q+2\textrm{diag}(\lambda)-2Z+\sum_{t=1}^m \gamma_t A_t\succeq 0,\\
			&\sigma\in\mathbb{R},~\alpha\in \mathbb{R}^n_+,~\beta\in \mathbb{R}^n_+,~\tau\in \mathbb{R}^m,\zeta\in \mathbb{R}^m_+,~M,N,R,S\in\mathbb{X}_+^n,\\
			&\lambda\in\mathbb{R}^n,~\gamma\in\mathbb{R}^m_+,~Z\in\mathbb{X}^n.
		\end{aligned}
	\end{equation}
	
	From any feasible solution $(\sigma,\alpha,\beta,\tau,\zeta,M,N,R,S,\lambda,\gamma,Z)$ of \eqref{QNR-dual3}, we can construct another solution $(\sigma,\alpha^\prime,\beta^\prime,\tau^\prime,\zeta,M,N,R,S,\lambda^\prime,\gamma^\prime,Z^\prime)$ by setting $\alpha^\prime=0$, $\beta^\prime=0$, $\tau^\prime=0$, $\gamma^\prime=\gamma-\tau$, $\lambda^\prime=\lambda+\beta$, and $Z^\prime=M+N-R-S$. We can check that the constructed solution is also feasible to \eqref{QNR-dual3}, and has the same objective value as the original solution. Thus, we may fix the vectors $\alpha$, $\beta$, and $\tau$ to zero to show that \eqref{QNR-dual3} is equivalent to \eqref{QNR-dual1}, and the first equation in  \eqref{eq_condition} is reduced to $Z=M+N-R-S$.
\end{proof}

\end{document}
